\documentclass[10pt]{article}

\usepackage[margin=1in]{geometry}
\usepackage{amsmath}
\usepackage{amssymb}
\usepackage{amsthm}
\usepackage{mathtools}
\usepackage{bm}

\usepackage{color}
\usepackage{colortbl}
\definecolor{deepblue}{rgb}{0,0,0.5}
\definecolor{deepred}{rgb}{0.6,0,0}
\definecolor{deepgreen}{rgb}{0,0.5,0}
\definecolor{gray}{rgb}{0.7,0.7,0.7}

\usepackage{hyperref}
\hypersetup{
  colorlinks   = true, %Colours links instead of ugly boxes
  urlcolor     = black, %Colour for external hyperlinks
  linkcolor    = blue, %Colour of internal links
  citecolor    = blue  %Colour of citations
}

%%%%%%%%%%%%%%%%%%%%%%%%%%%%%%%%%%%%%%%%%%%%%%%%%%%%%%%%%%%%%%%%%%%%%%%%%%%%%%%%

\theoremstyle{definition}
\newtheorem{problem}{Problem}
\newtheorem{defn}{Definition}
\newtheorem{note}{Note}
\newtheorem{refr}{References}
\newcommand{\E}{\mathbb E}
\newcommand{\R}{\mathbb R}
\DeclareMathOperator{\nnz}{nnz}
\DeclareMathOperator{\determinant}{det}
\DeclareMathOperator{\Var}{Var}
\DeclareMathOperator{\rank}{rank}
\DeclareMathOperator*{\argmin}{arg\,min}
\DeclareMathOperator*{\argmax}{arg\,max}

\newcommand{\I}{\mathbf I}
\newcommand{\Q}{\mathbf Q}
\newcommand{\p}{\mathbf P}
\newcommand{\pb}{\bar {\p}}
\newcommand{\pbb}{\bar {\pb}}
\newcommand{\pr}{\bm \pi}

\newcommand{\trans}[1]{{#1}^{T}}
\newcommand{\loss}{\ell}
\newcommand{\w}{\mathbf w}
\newcommand{\wstar}{{\w}^{*}}
\newcommand{\x}{\mathbf x}
\newcommand{\y}{\mathbf y}
\newcommand{\lone}[1]{{\lVert {#1} \rVert}_1}
\newcommand{\ltwo}[1]{{\lVert {#1} \rVert}_2}
\newcommand{\lp}[1]{{\lVert {#1} \rVert}_p}
\newcommand{\linf}[1]{{\lVert {#1} \rVert}_\infty}
\newcommand{\lF}[1]{{\lVert {#1} \rVert}_F}

\newcommand{\ignore}[1]{}

%%%%%%%%%%%%%%%%%%%%%%%%%%%%%%%%%%%%%%%%%%%%%%%%%%%%%%%%%%%%%%%%%%%%%%%%%%%%%%%%

\begin{document}


\begin{center}
{
\Huge
Chapter 1: The Learning Problem
}

%\vspace{0.15in}
%Due: Sunday, 6 Sep 2020 at midnight
\end{center}

\begin{center}
\includegraphics[width=\textwidth]{dilbert}
\end{center}

\ignore{
Exercise 1.1, 1.2, 1.3, 1.5, 1.6

Memorize the Hoefding Inequality

Exercise 1.11?

Exercise 1.12, 1.13?

Problem 1.2?
}


\section{Overview}

The first chapter of the book introduces
\begin{enumerate}
    \item basic notation for machine learning, and
    \item a learning model called the perceptron.
\end{enumerate}
We analyze the runtime properties of the perceptron in this chapter,
and we'll analyze the statistical properties of the algorithm in Chapter 2.
Then in Chapters 3-4, we'll look at more complicated learning algorithms.
%Most students find chapter 2 to be the most difficult content to understand.

\begin{note}
I have personally purchased for you all a physical copy of this textbook to emphasize the importance of reading.
You will not do well in this course if you do not read the textbook.
You should read each chapter at least once before we cover it in class,
and once after we cover it in class.
These notes packets follow the textbook very closely,
and I have annotated most of the problems with the corresponding sections in the textbook where you can find their solutions.

This textbook likely contains material that is much harder than the material you have encountered in previous math courses.
I have selected it because it is the only undergraduate-focused textbook that actually covers the math behind why machine learning works.
All other textbooks are focused on graduate students and significantly more ``dense.''
That said, this textbook is still quite dense.
A reading pace of 5 pages per hour would be fast.
You may have to reread certain sections several times, and that is normal.
\end{note}

\begin{note}
The author provides lecture videos that go along with the book, which you may find helpful.
They are located at
\begin{center}
\url{https://www.youtube.com/watch?v=mbyG85GZ0PI}.
\end{center}
%The author's general approach is a more ``traditional computer science'' approach to this material.
%That is, he presents a variety
The author has a very different emphasis than I do.
He focuses more on algorithm implementation details and proofs of theorems,
I focus more on understanding when to use each algorithm and when to apply the theorems.
\end{note}

\begin{note}
The first part of this course (on computational linear algebra and pagerank) is all background knowledge required for this textbook.
The details of the pagerank algorithm will not be needed for this book,
but an understanding of (1) the runtimes of linear algebra computations and (2) basic concepts like linear independence, rank, and eigenvectors are all assumed.
Part of the reason we began with the study of pagerank is that it helped you review those concepts in preparation for this book.
The textbook will add probability/statistics onto this background knowledge,
but most students find the probability/statistics to be much easier than the linear algebra.
(And probability and statistics are not prerequisites for this course.)
\end{note}


\newpage
\section{Section 1.1: Problem Setup}

\begin{problem}
    You are a bank and need an algorithm for determining whether to approve a credit loan application.
    Describe the learning framework notation for this problem.
    (See Section 1.1.1 and Figure 1.2 in the textbook.)

    NOTE: You should be able to do this for other applications as well.
    See Exercise 1.1 for a list of alternative applications to practice with.
\end{problem}

\newpage
\begin{problem}
    This problem explores the perceptron hypothesis class.
    \begin{enumerate}
    \item
    Define the Perceptron hypothesis class.
    (See Section 1.1.2 in the textbook.)
            \vspace{4in}

    \item
    Show ``good'' and ``bad'' examples of hypotheses in the Perceptron hypothesis class in $d=2$ dimensions.
    (See the discussion in Figure 1.3 and 3.1 in the textbook.)

    %\item
    %Show a sample data set for which the Perceptron hypothesis class will fail.
    %(See the discussion in Figure 3.1 in the textbook; note that this is not in the first chapter.)

            \vspace{4in}

    \item
    Describe which features for the bank loan problem are likely to have high weights, and which features are likely to have low weights on a ``good'' perceptron hypothesis.

    NOTE:
        This is closely related to Exercise 1.2 in the book, which you should be able to complete on your own.

        \newpage
    \end{enumerate}
\end{problem}

%\begin{problem}
    %Complete Exercise 1.2.
%\end{problem}


\newpage
\begin{problem}
    The Perceptron Learning Algorithm (PLA) is a procedure for selecting a particular hypothesis from the Perceptron hypothesis class.
    This problem is designed to help you understand the PLA.
    \begin{enumerate}
        \item State the PLA update rule, as found in Eq (1.3).
            \vspace{4in}
        \item Exercise 1.3 in the textbook is designed to help you get some intuition for why the PLA update rule is a good rule.
            It has the following parts:
            \begin{enumerate}
                \item Show that $y(t)\trans\w(t)\x(t) < 0$.
                    \vspace{4in}
                \item Show that $y(t)\trans\w(t+1)\x(t) > y(t)\trans\w(t)\x(t)$.
                    \vspace{4in}
                \item As far as classifying $\x(t)$ is concerned, argue that the move from $\w(t)$ to $\w(t+1)$ is a move `in the right direction'.

                    NOTE: See the Figure on the bottom of page 7 for a visual description of these results.
                    \vspace{4in}
            \end{enumerate}
            \newpage
        \item 
            Problem 1.3 (at the end of the chapter) proves that the total number of iterations $t$ required for the PLA to converge to an optimal solution is bounded by
%Problem 1.3
%Prove that the PLA eventually converges to a linear separator for seperable data.
%The following steps will guide you through the proof.
            \begin{equation}
                t \le \frac{R^2 \ltwo{\wstar}^2}{\rho^2},
            \end{equation}
            where $\wstar$ is an optimal hyperplane,
    \begin{equation}
        \rho = \min_{1\le n\le N} y_n(\trans\wstar\x_n)
    %\end{equation}
            \qquad\text{and}\qquad
    %\begin{equation}
        R = \max_{1\le n\le N} \ltwo{\x_n}.
    \end{equation}
            In this class, we won't complete that problem and prove this bound;
            instead, we will focus on understanding the consequences of this bound.
            \begin{enumerate}
                \item
                    Illustrate the quantities above.
    %Illustrate the optimal hyperplane $\wstar$,
    %and the quantities
    %See Problem 1.3 at the end of the chapter for a discussion of these formulas.
    %If you're having trouble visualizing these quantities,
    %complete Problem 1.2.

                    \vspace{4in}
            %\newpage
                %\item
                    %How are $\rho$ and $R$ likely to change as the number of feature dimensions change?
                    %\vspace{4in}
                \item
            What is the overall runtime of running the PLA to convergence?
            \end{enumerate}

            \vspace{4in}
            %\newpage
        \item
            Describe some of the differences between our runtime analysis for the PLA and the Pagerank power method algorithm.

        \end{enumerate}
\end{problem}
\newpage
\begin{problem}
    Based on the analysis above, draw two data sets that we expect PLA to have very different runtimes for, and explain which data set will be faster and which will be slower.
\end{problem}

\newpage
\begin{problem}
    %Preprocessing data points is a common
    In order to speed up the runtime of PLA,
    Alice is proposing several different variants of the algorithm.
    %For each variant below, explain why the runtime i
    For each variant below, identify the conditions (if any) under which it will have a faster runtime.
    \begin{enumerate}
        \item
            % FIXME:
            % Should we also add correctness changes to the PLA?
            % e.g. add some modifications that DO make it faster, but also make the results different?
            Alice proposes ``shrinking'' the data points by a factor of $\alpha$ so that $R$ will be smaller.
            That is, she first creates new data points defined by
            \begin{equation}
                \x'_i = \frac{\x_i}{\alpha}
            \end{equation}
            and then runs the PLA on this new data set.
            %Will the runtime of PLA on this modified data set be any faster?
            %Her idea is that the new radius of the data set will be $R'=R/\alpha$, and so the runtime above will be reduced by a factor of $\alpha$.
            %Why doesn't this transformation work?
            \newpage

        \item
            %In order to speed up the runtime of PLA,
            Alice proposes ``centering'' the data points so that $R$ will be smaller.
            That is, she first creates new data points defined by
            \begin{equation}
                \x'_i = \x_i - \bar\x
            \end{equation}
            where
            \begin{equation}
                \bar\x = \frac1N \sum_{i=1}^N \x_i,
            \end{equation}
            and then runs the PLA on this new data set.
            %Will the runtime of PLA on this modified data set be any faster?
        \newpage

        \item
            Alice notices that there is no unique value for $\wstar$ that defines the optimal separating hyperplane.
            In particular, for every $\wstar$ that defines a separating hyperplane and every positive real number $\alpha$,
            $\alpha\wstar$ also defines a separating hyperplane.
            By setting $\alpha$ to be small and using $\alpha\wstar$ as the separating hyperplane in the runtime formulas,
            Alice hopes to achieve a faster runtime.
            Why won't this work?
    \end{enumerate}
\end{problem}

\ignore{
\newpage
\section{Section 1.1.3: Learning vs Design}

\begin{problem}
You should complete Exercise 1.5 to get practice identifying when to use a learning approach to a problem, and when to use a ``design'' approach to a problem.
\end{problem}

\section{Section 1.2: Types of Learning}

The textbook identifies 3 types of learning: supervised, reinforcement, and unsupervised.
Pagerank is an example of an unsupervised algorithm (because we do not have ``labels'' about which nodes in the graph are the most important).
The focus of the rest of this class will be on supervised methods.

\begin{problem}
    You need to be able to identify examples of the different learning types.
    Exercise 1.6 will give you practice with this.
\end{problem}

\section{Section 1.3: Is Learning Feasible?}

\begin{problem}
    Complete Exercise 1.7.
\end{problem}

\begin{problem}
    Define the
    \begin{enumerate}
        \item Hoeffding Inequality (page 19, Eq 1.4)
        \item in-sample error (page 21)
        \item out-of-sample error (page 21)
        \item Eq 1.5
        \item union bound (page 24)
        \item Eq 1.6
    \end{enumerate}
\end{problem}

\begin{problem}
    Complete Exercise 1.11.
\end{problem}

\begin{problem}
    Complete Exercise 1.12.
\end{problem}

\section{Section 1.4: Error and Noise}

\begin{problem}
    Complete Exercise 1.13
\end{problem}
}

\ignore{
Problem 1.12

Chapter 2

Definition 2.1 dichotomies

Definition 2.2 growth function

Definition 2.3 break point

Definition 2.5 VC dimension

Exercise 2.4: VC dimension of PLA

Theorem 2.5

Exercise 2.5 substitute values into VC theorem

Theorem 2.4 + Eq (2.9) + substitute into Theorem 2.5

Figure 2.3: Error 

Exercise 2.6

Exercise 2.7 (do one in class)?

Section 2.3: bias/variance

Exercise 2.8 (proofs of gbar)

Problem 2.1, 2.2, 2.3, 2.7, 2.8, 2.10?, 2.11, 2.12, 2.13?, 2.14?, 2.15, 2.16, 2.17, 2.18, 
}

\end{document}


